\subsection{Uniform positive characters from small Bernoulli numbers}
\label{sec:positive-iwasawa-components}

\begin{proposition}[K17-POSITIVE-COMPONENTS]
\label{prop:positive-iwasawa-components}
For every odd prime $p$, every $n\geq0$, and every even integer
$a$ with $2\leq a\leq32$, one has
\[
                      A_n^{[a]}=0,
                      \qquad X^{[a]}=0.
 \tag{PC.1}
\]
The notation is that of Proposition~\ref{prop:componentwise-iwasawa-control}.
For $p\geq5$, the actual $p$-ramified component $M^{[1-a]}$
is integrally free of rank one over $\Lambda$.
For every positive $k$ satisfying
\[
       k\equiv-j\pmod{(p-1)/2}
                  \quad\text{for some }1\leq j\leq16,
\]
one also has $K_{4k}(\mathbf Z)\{p\}=0$.
These assertions use the established small-prime Vandiver verification
and the proved reflection lemma; they do not use Alp\"oge's theorem.
\end{proposition}

\begin{proof}
We first certify the complete prime support of the required Bernoulli
numerators. Use $B_0=1$, $B_1=-1/2$, and the exact recurrence
\[
 B_r=-\frac1{r+1}\sum_{i=0}^{r-1}\binom{r+1}{i}B_i
                         \qquad(r\geq1).
 \tag{PC.2}
\]
It follows directly by equating coefficients in
$(\exp(t)-1)\sum_{r\geq0}B_rt^r/r!=t$ and uniquely determines
every rational $B_r$. The resulting absolute numerator factors are
\[
\begin{array}{c|l@{\qquad}c|l}
a&|\operatorname{num}B_a|&a&|\operatorname{num}B_a|\\ \hline
2&1&18&43867\\
4&1&20&283\cdot617\\
6&1&22&11\cdot131\cdot593\\
8&1&24&103\cdot2294797\\
10&5&26&13\cdot657931\\
12&691&28&7\cdot9349\cdot362903\\
14&7&30&5\cdot1721\cdot1001259881\\
16&3617&32&37\cdot683\cdot305065927.
\end{array}
 \tag{PC.3}
\]
The complete standalone certificate
\path{work/bernoulli_positive_components_certificate.py}
computes (PC.2) using exact rational arithmetic, verifies the
factor products in (PC.3), and proves each displayed factor
prime by trial division through its integer square root.
The largest factor is $1001259881<2^{31}$. It also records every
denominator and checks the exact Faulhaber identity at three
integer arguments as a separate consistency check. No probable-prime
or numerical Bernoulli computation enters this certificate.

Suppose first that $p>2^{31}$. For the specified $a$, put
$e=p-1-a$. Then $e$ is even and $2\leq e\leq p-3$.
The denominator of $B_a$ is prime to $p$ (also directly visible
in the recorded rational values), and (PC.3) gives
$B_a\not\equiv0\pmod p$. Since $p-1-e=a$, the exact
Herbrand-and-reflection implication in
Lemma~\ref{lem:isolated-reflection} gives
\[
                         A_0^{[-e]}=A_0^{[a]}=0.
\]
This applies the lemma at the stated large value of $e$;
it does not apply a finite-weight Hecke calculation at that value.

For every odd $p<2^{31}$, the cited Hart--Harvey--Ong
verification gives Vandiver, hence every even component of
$A_0$ is zero. This includes all small-prime coincidences of
characters and all numerator or denominator divisors in (PC.3).
There is no prime equal to $2^{31}$. Thus $A_0^{[a]}=0$ for
all the asserted $p,a$. Proposition~\ref{prop:componentwise-iwasawa-control}
now proves (PC.1) at every finite layer and in the norm limit.

For $p\geq5$, its integral freeness implication gives
$M^{[1-a]}\simeq\Lambda$ whenever $[a]\ne[0],[2]$.
For the exceptional characters $[a]=[0]$ or $[2]$, the
corresponding odd characters are $[1]$ and $[-1]$; the complete
Coleman calculation in Corollary~\ref{cor:odd-companion-exclusion}
gives the same integral rank-one freeness. This proves the
assertion for every residue, including repetitions.

Finally let $a=2j$ and suppose
$k\equiv-j\pmod{(p-1)/2}$. Then
$-2k\equiv2j=a\pmod{p-1}$. The exact comparison (K2)
and $A_0^{[a]}=0$ give the asserted primary $K$-group
vanishing. This is a vanishing statement in residue classes,
not an assertion of periodicity of the group orders.
\end{proof}

\begin{corollary}[K17-TWO-SIDED-COMPONENT-RANGE]
\label{cor:two-sided-component-range}
With the stipulated Alp\"oge theorem, for every odd prime $p$
and positive integer $k$ whose residue modulo $(p-1)/2$
has a representative in the interval
\[
                              -16\leq j\leq9,
 \tag{PC.4}
\]
one has $K_{4k}(\mathbf Z)\{p\}=0$ and
$X^{[-2k]}=0$. For $p\geq5$ the corresponding
$M^{[2k+1]}$ is integrally free of rank one. These are precisely
the combined proved negative and positive character ranges;
the remaining characters are not asserted to vanish.
\end{corollary}

\begin{proof}
For $-16\leq j\leq-1$ use
Proposition~\ref{prop:positive-iwasawa-components}; for
$0\leq j\leq9$ use Corollary~\ref{cor:odd-companion-exclusion}
and Corollary~\ref{cor:nine-degrees-period}. The exponent
$-2k$ and its reflected exponent $2k+1$ depend only on the
specified residue of $k$. This proves every stated assertion.
\end{proof}
